% abstract.tex

\begin{abstract}
Retrasar la sincronización entre nodos en un entrenamiento distribuido reduce la frecuencia de comunicación pero induce divergencia entre los pesos (\emph{client drift}). Este trabajo evalúa el impacto de variar las épocas offline ($E$) utilizando el sistema de entrenamiento O.N.O. bajo las topologías \emph{Parameter Server} sincrónico y \emph{All-Reduce (ring)}. Los experimentos se ejecutaron sobre un clúster de contenedores Docker con emulación de red mediante Pumba. Entrenando LeNet-5 sobre MNIST FULL, los resultados muestran que incrementar $E$ no aporta reducciones significativas en los tiempos totales debido a la contención del hardware local, mientras que degrada la \emph{accuracy} final e incrementa la inestabilidad de la pérdida durante el entrenamiento.
\end{abstract}

\begin{IEEEkeywords}
Épocas Offline, Parameter Server, All-Reduce, O.N.O., Client Drift, Pumba.
\end{IEEEkeywords}
